\begin{claim}{D.2.10}{Proved}{Stable sphere annihilation implies finite-target Miller nullity}
Put
\[
 \mathcal A_{\mathbb S}
 =\{E\in\Ho(\Sp):\operatorname{F}(E,\mathbb S)\simeq0\},
\]
where $\operatorname{F}$ denotes the derived function spectrum, so the
condition includes every shift. If a pointed space $Y$ satisfies
$\Sigma^\infty Y\in\mathcal A_{\mathbb S}$, then, for every finite simply
connected pointed complex $K$ and every prime $p$,
\[
 \Map_*(Y,K)\simeq *,\qquad \Map_*(Y,K_{(p)})\simeq *.
\]
Consequently
\[
 R(\mathcal A_{\mathbb S})\subseteq W_S,
 \qquad R(\mathcal I_p)\subseteq W_S,
 \quad \mathcal I_p=\loc\bigl(\Sigma^\infty K(\FF_p,2)\bigr).
\]
The second inclusion uses D.2.6. Thus finite simply connected targets,
including the sphere tests defining $W_S$, cannot separate the
zero-finite-part candidate from its stable closure. This does not decide
whether that candidate is fixed by $RT$.
\end{claim}
\begin{proof}
Write $\mathcal C=\operatorname{coloc}(\mathbb S)$ for the smallest full
stable subcategory containing $\mathbb S$ and closed under retracts and
small homotopy limits. For each fixed $E$, the kernel of
$\operatorname{F}(E,-)$ has these closure properties. We will show that
the spectra underlying the Goodwillie layers of $K$ and $K_{(p)}$ lie
in $\mathcal C$.

\emph{Finite-group classifying spectra.}
For a finite $p$-group $P$, let $\mathbb S_P$ be the genuine equivariant
sphere. The $p$-completed Segal theorem gives
\begin{equation}\label{eq:D-stable-Miller-Segal}
 ((\mathbb S_P)^P)^{\wedge}_p
 \simeq \operatorname{F}(\Sigma^\infty BP_+,\mathbb S^{\wedge}_p).
\end{equation}
The precise comparison is \cite[Theorem~1.1]{FGHBergsakerRognes2023}
with $B=\mathbb S$; its introduction explicitly identifies $B^{\wedge P}$
with $\mathbb S_P$. Completion commutes here with the function spectrum:
use $X^{\wedge}_p=\holim_a X/p^a$, exactness, and dualizability of the
finite Moore spectra. The tom Dieck splitting
\cite[Corollary~2.3 and Section~3]{FGHCarlsson1992} is
\[
 (\mathbb S_P)^P\simeq
 \bigvee_{(H\leq P)}\Sigma^\infty B(N_P(H)/H)_+.
\]
The $H=1$ summand and \eqref{eq:D-stable-Miller-Segal} show that
$(\Sigma^\infty BP_+)^{\wedge}_p$ is a retract of a homotopy limit of
$\mathbb S^{\wedge}_p$. Since
$\mathbb S^{\wedge}_p=\holim_a\mathbb S/p^a$ belongs to $\mathcal C$,
so does this retract.

Let $G$ be any finite group and choose a Sylow $p$-subgroup $P$.
The stable transfer of the finite covering $BP\to BG$ gives
\[
 \Sigma^\infty BG_+\xrightarrow{\operatorname{tr}}
 \Sigma^\infty BP_+\xrightarrow{j}\Sigma^\infty BG_+.
\]
On ordinary homology, $j\operatorname{tr}$ is multiplication by
$[G:P]$: a singular simplex has exactly that many lifts. Thus this
composite becomes an equivalence after $p$-completion, and
$(\Sigma^\infty BG_+)^{\wedge}_p$ is a retract of
$(\Sigma^\infty BP_+)^{\wedge}_p$. This uses the homology formula for
the transfer: its cofiber is bounded below, and vanishing mod-$p$
homology implies that its smash product with $\mathbb S/p$ vanishes
by the stable Hurewicz theorem. It therefore has zero $p$-completion.
No equality of the stable composite with an integer multiple of the
identity is needed. The stable transfer construction is recalled
in \cite[Section~2(B)]{FGHCarlsson1992}.

The reduced spectrum $\Sigma^\infty BG$ has finite
$|G|$-primary homotopy groups in each degree. Indeed its connective
Atiyah--Hirzebruch spectral sequence has terms
$\widetilde H_s(BG;\pi_t\mathbb S)$, with $s\geq1$, $t\geq0$,
and finitely many terms in each total degree. These groups are finite
and annihilated by $|G|$: use the bar resolution, which is finite in
each degree, and the ordinary transfer to the trivial group.
Consequently the natural map
\[
 \Sigma^\infty BG\longrightarrow
 \prod_{p\mid |G|}(\Sigma^\infty BG)^{\wedge}_p
\]
is an equivalence on every homotopy group. This finite product, and
then $\Sigma^\infty BG_+\simeq\mathbb S\vee\Sigma^\infty BG$, belong
to $\mathcal C$.

\emph{Localization of the orbit cells.}
We also need $(\Sigma^\infty BG_+)_{(p)}\in\mathcal C$.
First the cofiber of
\[
 \mathbb S\longrightarrow\prod_q\mathbb S^{\wedge}_q
\]
is $HV$, where $V=(\prod_q\ZZ_q)/\ZZ$. This follows on homotopy
groups: the map is an isomorphism in positive degrees and injective
in degree zero. The group $V$ is a nonzero rational vector space.
For a prime $q$, multiplication by $q$ is surjective by subtracting
an integer with the prescribed residue in the $q$-coordinate; it is
injective because a rational number with denominator $q$ lying in
every $\ZZ_r$ must be an integer. Hence $H\QQ$ is a retract of $HV$
and lies in $\mathcal C$. Every rational vector space embeds as a
direct summand in a product of copies of $\QQ$, by choosing a basis
and splitting the inclusion of its direct sum into the product.
Applying $H$ proves that $HU\in\mathcal C$ for every rational vector
space $U$. This uses products and retracts, with no coproduct-closure
assumption on $\mathcal C$.

The arithmetic pullback
\[
 \mathbb S_{(p)}\simeq
 \mathbb S^{\wedge}_p\mathop{\times}^{h}_{H\QQ_p}H\QQ
\]
therefore lies in $\mathcal C$. One can check this pullback directly:
in positive degrees it has the $p$-primary stable stems, while
$\ZZ_p\cap\QQ=\ZZ_{(p)}$ and $\ZZ_p+\QQ=\QQ_p$ give degree zero
and the vanishing of degree $-1$. Since the reduced classifying
spectrum has degreewise finite homotopy groups,
$(\Sigma^\infty BG)_{(p)}\simeq(\Sigma^\infty BG)^{\wedge}_p$.
Adding $\mathbb S_{(p)}$ proves the required localized assertion.

It follows that if $F$ is a genuine finite $G$-spectrum, then both
$F_{hG}$ and $(F_{hG})_{(p)}$ lie in $\mathcal C$. Indeed finite
genuine spectra are obtained by finite cofiber constructions and
retracts from the cells $G/H_+\wedge\mathbb S^d$, whose homotopy
orbits are $\Sigma^d\Sigma^\infty BH_+$. This assertion is only
needed for genuinely finite equivariant models; it does not assume
that an arbitrary action on an underlying finite spectrum has such
a model.

\emph{The identity tower.}
The identity functor has layers
\[
 D_k I(X)\simeq
 \Omega^\infty\bigl(\partial_k I\wedge
                    (\Sigma^\infty X)^{\wedge k}\bigr)_{h\Sigma_k}.
\]
By \cite[Theorem~2.1]{EH5AroneMahowald1999}, $\partial_k I$ is the
Spanier--Whitehead dual of a finite partition complex with its
specified $\Sigma_k$-action. This is the restriction of a genuine
finite $\Sigma_k$-spectrum. A finite simplicial model of $K$ gives a
finite equivariant CW model of $K^{\wedge k}$ after subdivision;
duality and smash products preserve genuine finiteness. Thus the
spectrum in the displayed layer for $X=K$ belongs to $\mathcal C$.
For $X=K_{(p)}$, use the natural equivalence
$\Sigma^\infty K_{(p)}\simeq(\Sigma^\infty K)_{(p)}$ and the
localized orbit assertion just proved. The layer formula extends
from finite complexes to all spaces because the identity and its
layers preserve filtered homotopy colimits.

Put $E=\Sigma^\infty Y$. Our hypothesis implies
$\operatorname{F}(E,V)=0$ for every $V\in\mathcal C$.
Adjunction therefore makes $\Map_*(Y,D_kI(X))$ contractible in both
cases. Starting from $P_0I(X)=*$, the fibre sequences of the tower
give $\Map_*(Y,P_kI(X))\simeq *$ for all $k$. The identity is
$1$-analytic \cite[Example~4.3]{FGHGoodwillie1992}; hence its tower
converges on the simply connected spaces $K$ and $K_{(p)}$
\cite[Theorem~1.13]{FGHGoodwillie2003}. Mapping out of $Y$ preserves
this homotopy limit, proving both mapping-space assertions without
any compactness assumption on $Y$.

Finally, for $f$ in $R(\mathcal A_{\mathbb S})$, apply the assertion
to its pointed homotopy cofiber $C(f)$ and to $K=S^n$, $n\geq2$,
at $p=2$. This is exactly the cofiber condition defining $W_S$.
There is no component exception: a nonbasepoint component of a
pointed space makes $\mathbb S$ a retract of its reduced suspension
spectrum, so $\Sigma^\infty C(f)\in\mathcal A_{\mathbb S}$ already
forces $C(f)$ to be connected. The class $\mathcal A_{\mathbb S}$
is localizing, since $\operatorname{F}(-,\mathbb S)$ takes
homotopy colimits to homotopy limits. D.2.6 puts its generator
$\Sigma^\infty K(\FF_p,2)$ in this class, so
$\mathcal I_p\subseteq\mathcal A_{\mathbb S}$ and the last inclusion
follows.
\end{proof}
